\chapter{Обзор платформы и языка Java}
\label{ch:01}

Курс начинается с ответа на вопрос, как устроена Java. Это не только язык, но и
платформа: компилятор, виртуальная машина, стандартная библиотека и инструменты
разработчика; не понимая, как эти части связаны, правила языка приходится
заучивать наизусть. Прочитав главу, вы проследите путь кода от файла до строки
на экране, научитесь различать JDK, JRE и JVM, узнаете, кто загружает классы в
память, где живут объекты и почему в Java не нужно освобождать память вручную.
Завершают главу разбор системы типов и операторов~--- фундамента всех
последующих глав.

\section{Java как платформа: байт-код и переносимость}

Java~--- объектно-ориентированный язык: код организован вокруг классов и
объектов, как дом из кубиков LEGO собирается из отдельных деталей, а не
выливается одним куском бетона. В отличие от C++, где функцию можно написать
«просто так», в Java любой код обязан находиться внутри класса: для программы в
три строки это избыточно, но на большом проекте только такая дисциплина
удерживает код от превращения в кашу.

Вторая особенность важнее~--- платформенная независимость. Компилятор C или C++
превращает исходный код в инструкции конкретного процессора, поэтому программа,
собранная для Windows на x86, не запустится ни на macOS, ни на телефоне с ARM.
Java превращает код в \term{байт-код} (\eng{bytecode})~--- промежуточный формат
для воображаемого процессора, который выполняет \term{виртуальная машина Java}
(\eng{Java Virtual Machine, JVM}). JVM написана для всех распространённых
систем, и один и тот же байт-код выполняется любой из них: программу переносят
с Windows на Linux готовым файлом \texttt{.class}, без пересборки. Отсюда девиз
языка~--- \eng{«Write once~--- run anywhere»}.

Путь от текста программы до результата показан на рис.~\ref{fig:01-1}; обратите
внимание, что компилятор и виртуальная машина~--- две разные программы,
запускаемые двумя разными командами.

\begin{figure}[htbp]\centering
\begin{tikzpicture}[font=\small,node distance=6mm]
\node[boxw=23mm](a){\texttt{Main.java}\\[1pt]\scriptsize исходный код};
\node[boxw=17mm,right=of a](b){\texttt{javac}\\[1pt]\scriptsize компилятор};
\node[boxw=23mm,right=of b](c){\texttt{Main.class}\\[1pt]\scriptsize байт-код};
\node[boxw=17mm,right=of c](d){\texttt{java}\\[1pt]\scriptsize JVM};
\node[boxw=21mm,right=of d](e){результат\\[1pt]\scriptsize в консоли};
\draw[flowarrow](a)--(b);\draw[flowarrow](b)--(c);
\draw[flowarrow](c)--(d);\draw[flowarrow](d)--(e);
\end{tikzpicture}
\caption{Путь от исходного кода к работающей программе}\label{fig:01-1}
\end{figure}

Сама JVM ничего не знает о языке Java~--- она умеет только выполнять байт-код.
Поэтому в байт-код компилируются и Kotlin, Scala, Groovy, а их классы живут в
одном проекте с классами Java.

\section{Первая программа и комментарии}

Знакомство с языком принято начинать с программы, печатающей приветствие
(листинг~\ref{lst:01-1}). Наберите её и сохраните в файл \texttt{Main.java}~---
ниже мы разберём каждую строку.

\begin{listing}[H]
\begin{minted}{java}
package lecture.one.hello;

// Имя файла обязано совпадать с именем публичного класса:
// класс Main лежит в файле Main.java
public class Main {

    // main - точка входа: с этого метода JVM начинает работу
    public static void main(String[] args) {
        System.out.println("Hello world!");
    }
}
\end{minted}
\caption{Программа, печатающая приветствие}
\label{lst:01-1}
\end{listing}

Первая строка объявляет \term{пакет} (\eng{package})~--- логическую группу
классов; пакет удобно представлять как полку в шкафу: носки на одной полке,
рубашки на другой. Имя пакета пишется строчными буквами и повторяет структуру
каталогов: класс из пакета \texttt{lecture.one.hello} лежит в папке
\texttt{lecture/one/hello}. Дальше объявлен класс: \texttt{public}~---
модификатор доступа, \texttt{class}~--- ключевое слово, \texttt{Main}~---
выбранное вами имя, и здесь действует жёсткое правило: имя файла должно
совпадать с именем публичного класса.

Метод \texttt{main}~--- \term{точка входа}: именно его ищет JVM при запуске, и
заголовок обязан выглядеть в точности так. Слово \texttt{public} открывает JVM
доступ к методу извне; \texttt{static} означает, что метод принадлежит классу, а
не объекту, и вызвать его можно, не создавая экземпляра,~--- иначе вышел бы
замкнутый круг: чтобы создать объект, надо выполнить код, а чтобы выполнить код,
нужен объект; \texttt{void} говорит, что метод ничего не возвращает;
\texttt{String[] args}~--- массив аргументов командной строки. В теле метода
\texttt{System}~--- класс доступа к системным ресурсам, \texttt{out}~--- поток
вывода, \texttt{println} печатает строку и переводит курсор на новую.

Строки с двумя косыми чертами~--- \term{комментарии}: компилятор их игнорирует.
Видов комментариев три (листинг~\ref{lst:01-2}); третий отличается от второго
удвоенной звёздочкой, но именно он умеет превращаться в документацию.

\begin{listing}[H]
\begin{minted}{java}
// Однострочный комментарий действует до конца строки

/* Многострочный комментарий
   может занимать сколько угодно строк */

/** Документирующий комментарий (Javadoc): из таких
 *  комментариев инструмент javadoc собирает справку. */
\end{minted}
\caption{Три вида комментариев}
\label{lst:01-2}
\end{listing}

Хорошее правило: комментарий объясняет \emph{почему} сделано именно так, а не
\emph{что} делает код,~--- последнее должно быть понятно из имён. Пометка
\texttt{// увеличиваем i на 1} рядом с \texttt{i++} бесполезна, а
\texttt{// места нумеруются с~1, компенсируем сдвиг} несёт контекст, которого в
коде нет. Сам инструмент \texttt{javadoc} разбирается в главе~\ref{ch:10}.

\section{JDK, JRE и JVM. Инструменты разработчика}

Три аббревиатуры~--- JDK, JRE и JVM~--- путают чаще всего, хотя устроены они
просто: это три матрёшки, вложенные одна в другую (рис.~\ref{fig:01-2}).
Читайте схему изнутри наружу: каждая следующая рамка добавляет к предыдущей
ровно одну возможность.

\begin{figure}[htbp]\centering
\begin{tikzpicture}[font=\small,sub/.style={font=\scriptsize}]
\draw[rounded corners=3pt](0,0)rectangle(11,3.9);
\node at (5.5,3.4){\bfseries JDK~--- комплект разработчика};
\node[sub] at (5.5,3.0){+ инструменты: javac, jar, javadoc, jshell};
\draw[rounded corners=3pt](0.6,0.4)rectangle(10.4,2.6);
\node at (5.5,2.1){\bfseries JRE~--- среда выполнения};
\node[sub] at (5.5,1.7){+ библиотеки классов: String, ArrayList, File};
\draw[rounded corners=3pt](1.2,0.7)rectangle(9.8,1.3);
\node at (5.5,1.0){\bfseries JVM~--- выполняет байт-код};
\end{tikzpicture}
\caption{Вложенность JDK, JRE и JVM}\label{fig:01-2}
\end{figure}

Самая маленькая матрёшка~--- \term{JVM}: программа, которая читает байт-код,
объясняет процессору, что делать, управляет памятью и следит, чтобы загружаемый
код не лез в чужую. Средняя матрёшка~--- \term{JRE}
(\eng{Java Runtime Environment}): JVM плюс стандартная библиотека из тысяч
готовых классов. Самая большая~--- \term{JDK} (\eng{Java Development Kit}): JRE
плюс инструменты разработки. Правило простое: хотите запускать~--- поставьте
JRE, хотите писать~--- поставьте JDK, который включает JRE.

Инструменты JDK перечислены в табл.~\ref{tab:01-1}. Первые два~--- обязательный
минимум, остальные понадобятся в следующих главах.

\begin{table}[htbp]
\caption{Основные инструменты JDK}
\label{tab:01-1}
\centering
\begin{tabular}{L{2.7cm}L{6.3cm}L{5.7cm}}
\toprule
Инструмент & Назначение & Пример вызова \\
\midrule
\texttt{javac} & компилирует код в байт-код & \texttt{javac Main.java} \\
\texttt{java} & запускает готовую программу & \texttt{java Main} \\
\texttt{jar} & упаковывает классы в архив & \texttt{jar cf app.jar *.class} \\
\texttt{javap} & показывает содержимое class-файла & \texttt{javap -c Main} \\
\texttt{javadoc} & собирает справку из комментариев & \texttt{javadoc *.java} \\
\texttt{jshell} & интерактивная консоль & \texttt{jshell} \\
\bottomrule
\end{tabular}
\end{table}

Упомянутый в таблице \term{JAR-файл} (\eng{Java ARchive})~--- обычный ZIP-архив
со скомпилированными классами: вместо сотни файлов \texttt{.class} коллеге
передают один \texttt{.jar}. А \texttt{jshell} (появился в Java~9) позволяет
проверять гипотезы, не создавая ни файла, ни класса, ни метода \texttt{main}: в
сеансе из листинга~\ref{lst:01-3} результат каждого выражения печатается сразу.

\begin{listing}[H]
\begin{minted}{text}
jshell> int x = 5
x ==> 5

jshell> x * 2
$2 ==> 10
\end{minted}
\caption{Сеанс работы в jshell}
\label{lst:01-3}
\end{listing}

Команда \texttt{/vars} печатает объявленные переменные, \texttt{/exit} завершает
сеанс. Но помните: \texttt{jshell}~--- инструмент для экспериментов, настоящие
программы состоят из классов и файлов \texttt{.java}. Профессиональная
разработка ведётся в \term{интегрированной среде разработки} (\eng{IDE}),
которая берёт на себя автодополнение, запуск одной кнопкой, отладку и
рефакторинг; в курсе используется IntelliJ~IDEA, её бесплатной редакции хватает
на все задания пособия.

\section{Внутреннее устройство JVM}

Виртуальная машина состоит из трёх подсистем: загрузчика классов, областей
памяти и исполняющего движка (рис.~\ref{fig:01-3}). Стрелки означают обмен
данными, а не порядок запуска: подсистемы работают одновременно.

\begin{figure}[htbp]\centering
\begin{tikzpicture}[font=\scriptsize,
  b/.style={draw,rounded corners=2pt,align=center,text width=20mm,
    minimum height=10mm,inner sep=2pt},
  t/.style={font=\scriptsize\bfseries,anchor=south},
  g/.style={draw,rounded corners=3pt,inner sep=2mm}]
\node[b](ma){Method Area\\(Metaspace)};
\node[b,right=2mm of ma](hp){Heap\\(куча)};
\node[b,right=2mm of hp](sk){JVM Stack\\(стек)};
\node[b,right=2mm of sk](pc){PC Register};
\node[b,right=2mm of pc](nm){Native Method Stack};
\node[fit=(ma)(nm),inner sep=0](ru){};
\node[t](rp)at([yshift=1.5mm]ru.north){Runtime Data Areas~--- память};
\node[g,fit=(ru)(rp)](rt){};
\node[b,text width=48mm,above=11mm of rt](cl)
  {\textbf{Class Loader}\\находит, проверяет и загружает \texttt{.class}};
\node[b,text width=26mm,below=15mm of rt](jt){JIT-компилятор};
\node[b,text width=26mm,left=3mm of jt](ip){Интерпретатор};
\node[b,text width=26mm,right=3mm of jt](gc){Сборщик мусора};
\node[fit=(ip)(gc),inner sep=0](eu){};
\node[t](ep)at([yshift=1.5mm]eu.north){Execution Engine~--- движок};
\node[g,fit=(eu)(ep)](ee){};
\draw[Latex-Latex](cl)--(rt);\draw[Latex-Latex](rt)--(ee);
\end{tikzpicture}
\caption{Внутреннее устройство виртуальной машины Java}\label{fig:01-3}
\end{figure}

\subsection{Загрузчик классов}

Запуская программу, JVM не загружает все классы разом: она подтягивает их по
мере надобности, и занимается этим \term{загрузчик классов} (\eng{class loader})
в три этапа.

На этапе \term{загрузки} (\eng{loading}) загрузчик находит файл \texttt{.class},
читает его и создаёт в памяти объект \texttt{Class}, описывающий поля, методы и
предка. На этапе \term{связывания} (\eng{linking}) выполняются три действия:
\emph{проверка} (\eng{verification})~--- JVM убеждается, что байт-код корректен и
безопасен; \emph{подготовка} (\eng{preparation})~--- под статические поля
выделяется память со значениями по умолчанию; \emph{разрешение}
(\eng{resolution})~--- символические имена заменяются прямыми ссылками. На этапе
\term{инициализации} (\eng{initialization}) выполняются статические блоки.

Загрузчик в JVM не один, их три, и каждый отвечает за свою зону.
\term{Bootstrap Class Loader} загружает базовые классы модуля
\texttt{java.base}~--- \texttt{String}, \texttt{Object}, \texttt{System};
написан он на C/C++ и является частью самой JVM, поэтому в виде Java-объекта не
существует: вызов \texttt{getClassLoader()} у класса \texttt{String} возвращает
\texttt{null}. \term{Platform Class Loader} (до Java~9~--- Extension) отвечает
за остальные модули платформы вроде \texttt{java.sql}, а
\term{Application Class Loader}~--- за то, что лежит в \texttt{classpath}: ваш
код и библиотеки.

Работают загрузчики по принципу \term{делегирования снизу вверх}
(\eng{parent-first delegation}): получив запрос, загрузчик не ищет класс сам, а
передаёт запрос родителю, тот~--- своему родителю, и так до корневого. Только
если родитель вернулся ни с чем, загрузчик ищет класс сам.

У такого порядка три важных следствия. Безопасность: подменить системный класс
невозможно~--- даже если вы напишете собственный \texttt{java.lang.String},
запрос сначала дойдёт до Bootstrap, и загружен будет настоящий класс.
Уникальность: каждый класс загружается ровно одним загрузчиком, конфликты версий
исключены. Видимость: классы родителя видны потомкам, но не наоборот. Проверить
всё это можно за две строки (листинг~\ref{lst:01-4}): в первой строке вывода
окажется \texttt{null}, и это не ошибка.

\begin{listing}[H]
\begin{minted}{java}
// Bootstrap - часть JVM, поэтому здесь будет null
System.out.println(String.class.getClassLoader());

// класс Main из вашего Main.java -> Application Loader
System.out.println(Main.class.getClassLoader());
\end{minted}
\caption{Кто какой класс загрузил}
\label{lst:01-4}
\end{listing}

\subsection{Области памяти}

Различать области памяти необходимо: от того, где лежит значение, зависят и
время его жизни, и поведение оператора сравнения.

\term{Куча} (\eng{heap})~--- общая область, где живут все объекты: каждый раз,
когда выполняется оператор \texttt{new}, объект создаётся именно здесь; куча
одна на всю программу и доступна всем потокам. \term{Стек} (\eng{stack}) у
каждого потока свой, в нём хранятся локальные переменные, параметры методов и
адреса возврата: вызов метода добавляет в стек новый кадр, завершение метода
кадр удаляет. \term{Счётчик команд} (\eng{PC register}) указывает на текущую
инструкцию потока, \term{Metaspace} хранит метаданные классов, а \term{стек
нативных методов}~--- вызовы кода на C/C++.

\begin{vrezka}
\textbf{Важно.} До Java~8 область метаданных называлась PermGen и имела
фиксированный размер, из-за чего появлялась ошибка
\texttt{OutOfMemoryError: PermGen space}. С Java~8 её заменил Metaspace: он
лежит в нативной памяти операционной системы и растёт по мере надобности,
поэтому ограничивающий кучу параметр \texttt{-Xmx} на него не действует.
\end{vrezka}

\subsection{Исполняющий движок}

\term{Интерпретатор} читает инструкции байт-кода одну за другой и выполняет
каждую: это просто, но медленно~--- если цикл выполняется миллион раз,
интерпретировать его приходится миллион раз. Поэтому в работу вступает
\term{JIT-компилятор} (\eng{Just-In-Time}): он находит «горячие точки»
(\eng{hot spots})~--- участки кода, которые вызываются особенно часто,~---
компилирует их в машинный код процессора и складывает в кэш, откуда они и
выполняются при следующем вызове. Отсюда название самой распространённой
реализации виртуальной машины~--- HotSpot JVM, и практическое следствие:
прогретая программа работает в разы быстрее, чем сразу после запуска.

Третий компонент~--- \term{сборщик мусора} (\eng{garbage collector, GC}): он
находит в куче объекты, на которые не осталось ни одной ссылки, удаляет их и
возвращает память программе. Именно поэтому в Java нет ни оператора
освобождения памяти, ни двух связанных с ним классических ошибок: утечки из-за
забытого освобождения и обращения к уже освобождённой памяти.

\section{Взаимодействие с нативным кодом: JNI}

Иногда чистой Java недостаточно: нужно вызвать функцию операционной системы или
подключить готовую библиотеку на C++. Для этого существует \term{JNI}
(\eng{Java Native Interface})~--- мост между миром Java и миром нативного кода:
вы объявляете метод с ключевым словом \texttt{native} и не пишете его тело, а
реализацию готовите на C или C++. В листинге~\ref{lst:01-5} важны две детали~---
отсутствие тела у метода и статический блок, загружающий библиотеку.

\begin{listing}[H]
\begin{minted}{java}
public class HelloJNI {

    // Статический блок выполняется один раз при загрузке
    static {
        // Ищет HelloJNI.dll (.so, .dylib) в java.library.path
        System.loadLibrary("HelloJNI");
    }

    // Тела нет: реализация написана не на Java
    public native void sayHello();
}
\end{minted}
\caption{Класс с нативным методом}
\label{lst:01-5}
\end{listing}

Заголовочный файл для C генерирует сам компилятор (ключ \texttt{-h} у
\texttt{javac}), а реализацию собирают компилятором C в разделяемую библиотеку.

\begin{oshibka}
\textbf{Когда не стоит браться за JNI.} Если задача решается на чистой Java~---
решайте её на чистой Java. Нативный код тяжело отлаживать, ошибка в нём роняет
не метод, а всю виртуальную машину, утечки памяти сборщик мусора не подберёт, а
собирать библиотеку придётся под каждую операционную систему отдельно~--- то
есть главное достоинство Java, переносимость, теряется.
\end{oshibka}

\section{Типы данных и преобразования}

Все типы Java делятся на две группы. \term{Примитивный тип} хранит само
значение: записав \texttt{int x = 5}, вы получаете переменную, внутри которой
лежит число 5. \term{Ссылочный тип} хранит адрес объекта: записав
\texttt{String s = "Hello"}, вы получаете переменную, внутри которой лежит не
текст, а адрес того места в куче, где текст находится. Остальные различия
вытекают из этого одного (табл.~\ref{tab:01-2}); главная строка в ней~---
предпоследняя, именно она порождает больше всего ошибок у начинающих.

\begin{table}[htbp]
\caption{Примитивные и ссылочные типы}
\label{tab:01-2}
\centering
\begin{tabular}{L{3.6cm}L{5.3cm}L{5.3cm}}
\toprule
Критерий & Примитивные & Ссылочные \\
\midrule
Что хранит переменная & само значение & адрес объекта в памяти \\
Где размещаются данные & локальная переменная~--- в стеке, поле~--- в куче внутри объекта & объект~--- всегда в куче, ссылка~--- как примитив \\
Значение по умолчанию & \texttt{0}, \texttt{false}, \texttt{'\textbackslash u0000'} & \texttt{null} \\
Оператор \texttt{==} & сравнивает значения & сравнивает адреса \\
Может быть \texttt{null} & нет & да \\
\bottomrule
\end{tabular}
\end{table}

Статические поля, в отличие от обычных, хранятся не в куче, а в области
метаданных классов.

Примитивных типов ровно восемь (табл.~\ref{tab:01-3}), и набор их фиксирован:
\texttt{int} всегда занимает 32~бита, какими бы ни были процессор и система.
Значение по умолчанию из последнего столбца получают только поля классов;
локальную переменную компилятор обязывает инициализировать явно.

\begin{table}[htbp]
\caption{Примитивные типы данных Java}
\label{tab:01-3}
\centering
\begin{tabular}{L{1.8cm}L{1.6cm}L{6.4cm}L{2.9cm}}
\toprule
Тип & Размер & Диапазон значений & По умолчанию \\
\midrule
\texttt{byte} & 8~бит & от $-128$ до $127$ & \texttt{0} \\
\texttt{short} & 16~бит & от $-32\,768$ до $32\,767$ & \texttt{0} \\
\texttt{int} & 32~бита & примерно $\pm 2{,}1$~млрд & \texttt{0} \\
\texttt{long} & 64~бита & примерно $\pm 9{,}2 \cdot 10^{18}$ & \texttt{0L} \\
\texttt{float} & 32~бита & дробное, 6--7 значащих цифр & \texttt{0.0f} \\
\texttt{double} & 64~бита & дробное, около 15 значащих цифр & \texttt{0.0} \\
\texttt{char} & 16~бит & один символ Unicode & \texttt{'\textbackslash u0000'} \\
\texttt{boolean} & --- & \texttt{true} или \texttt{false} & \texttt{false} \\
\bottomrule
\end{tabular}
\end{table}

\begin{oshibka}
\textbf{Частая ошибка: переполнение.} При выходе за границу диапазона значение
не обрезается до максимума, а «заворачивается» на противоположный край: строки
\texttt{byte b = 127; b++;} дают не 128, а $-128$, а прибавление единицы к
\texttt{Integer.MAX\_VALUE} даёт \texttt{Integer.MIN\_VALUE}. Исключения при
этом не возникает~--- программа продолжает считать с неверным числом.
\end{oshibka}

\begin{oshibka}
\textbf{Частая ошибка: точность дробных чисел.} Типы \texttt{float} и
\texttt{double} хранят числа в двоичном формате IEEE~754, а десятичную дробь 0,1
в двоичном виде точно записать нельзя. Поэтому \texttt{0.1 + 0.2} даёт
\texttt{0.30000000000000004}, а сравнение \texttt{0.1 + 0.2 == 0.3} возвращает
\texttt{false}. Для денежных расчётов берите класс \texttt{BigDecimal}.
\end{oshibka}

Оператор \texttt{==} сравнивает содержимое переменных, а в ссылочной переменной
лежит адрес; значит, для объектов он отвечает на вопрос «это один и тот же
объект?», а не «одинаковое ли у них содержимое»~--- содержимое сравнивает метод
\texttt{equals}. Путаницу добавляет \term{пул строк} (\eng{String Pool}):
одинаковые строковые литералы Java хранит в единственном экземпляре, причём
начиная с Java~7 сам пул лежит внутри кучи~--- это не отдельная область памяти,
а таблица уникальных строк в куче. Все три случая собраны в
листинге~\ref{lst:01-6}.

\begin{listing}[H]
\begin{minted}{java}
String s1 = "Java";
String s2 = "Java";              // тот же объект из String Pool
String s3 = new String("Java");  // новый объект в куче

System.out.println(s1 == s2);       // true  - адрес один
System.out.println(s1 == s3);       // false - адреса разные
System.out.println(s1.equals(s3));  // true  - текст совпадает
\end{minted}
\caption{Сравнение строк: == и equals}
\label{lst:01-6}
\end{listing}

Запомните сразу и навсегда: примитивы сравнивают через \texttt{==}, объекты~---
через \texttt{equals}. Кроме классов, к ссылочным типам относятся интерфейсы,
массивы, перечисления и аннотации.

\subsection{Преобразование числовых типов}

\term{Расширяющее преобразование} (\eng{widening}) идёт по цепочке
\texttt{byte} $\to$ \texttt{short} $\to$ \texttt{int} $\to$ \texttt{long} $\to$
\texttt{float} $\to$ \texttt{double} и выполняется само собой: компилятор
разрешения не спрашивает~--- это как перелить стакан воды в ведро. Оговорка
одна: на переходах \texttt{int} $\to$ \texttt{float} и \texttt{long} $\to$
\texttt{float}/\texttt{double} диапазон растёт, но точность теряться может~---
очень большое целое попадает в дробный тип округлённым. \term{Сужающее
преобразование} (\eng{narrowing}) идёт в обратную сторону и может потерять
данные, поэтому требует явного приведения в круглых скобках: это попытка
перелить ведро в стакан, и компилятор заставляет письменно подтвердить, что
часть воды прольётся (листинг~\ref{lst:01-7}).

\begin{listing}[H]
\begin{minted}{java}
int i = 100;
long l = i;            // расширяющее: выполняется само

double x = 3.9;
int y = (int) x;       // сужающее: 3, дробная часть отброшена
byte b = (byte) 300;   // 44 - лишние биты обрезаны

char c = 'A';
int code = c + 1;           // 66: char продвинут до int
char next = (char) (c + 1); // 'B' - обратно нужно явно
\end{minted}
\caption{Расширяющее и сужающее преобразования}
\label{lst:01-7}
\end{listing}

Правило продвижения операндов таково: если хотя бы один операнд имеет тип
\texttt{byte}, \texttt{short} или \texttt{char}, оба продвигаются минимум до
\texttt{int}. Отсюда две неожиданности из листинга~\ref{lst:01-8}: сложение двух
\texttt{byte} даёт \texttt{int}, а составное присваивание \texttt{+=} молча
добавляет приведение, которого вы не писали.

\begin{listing}[H]
\begin{minted}{java}
byte n1 = 10, n2 = 20;
// byte sum = n1 + n2;  // ошибка: результат имеет тип int
int sum = n1 + n2;

byte value = 10;
value += 300;          // на деле: value = (byte)(value + 300)
System.out.println(value);   // 54, а не 310
\end{minted}
\caption{Продвижение операндов и составное присваивание}
\label{lst:01-8}
\end{listing}

Именно поэтому \texttt{value = value + 300} не компилируется, а
\texttt{value += 300} компилируется: во втором случае приведение
\texttt{(byte)} компилятор подставляет за вас.

Наконец, о записи самих чисел: Java понимает четыре системы счисления и
разрешает разделять разряды подчёркиванием (листинг~\ref{lst:01-9}). На значение
подчёркивания не влияют~--- они существуют для читателя.

\begin{listing}[H]
\begin{minted}{java}
int decimal = 42;        // десятичная запись
int hex     = 0x2A;      // шестнадцатеричная, префикс 0x
int octal   = 052;       // восьмеричная, префикс 0
int binary  = 0b101010;  // двоичная, префикс 0b
// все четыре переменные хранят одно и то же число 42

int  million = 1_000_000;   // подчёркивания - только для глаз
long big     = 123456789L;  // L обязателен для long
float pi     = 3.14f;       // f обязателен для float
\end{minted}
\caption{Числовые литералы}
\label{lst:01-9}
\end{listing}

\section{Операторы}

Операторы~--- знаки действий над данными. Ведут они себя в основном
по-школьному: \texttt{*}, \texttt{/} и \texttt{\%} выполняются раньше
\texttt{+} и \texttt{-}, а скобки перебивают любой приоритет. Арифметических
операторов пять (табл.~\ref{tab:01-4}); сюрприз готовит только деление.

\begin{table}[htbp]
\caption{Арифметические операторы}
\label{tab:01-4}
\centering
\begin{tabular}{L{2.3cm}L{6.6cm}L{5.1cm}}
\toprule
Оператор & Действие & Пример \\
\midrule
\texttt{+} & сложение & \texttt{10 + 3} $\to$ \texttt{13} \\
\texttt{-} & вычитание & \texttt{10 - 3} $\to$ \texttt{7} \\
\texttt{*} & умножение & \texttt{10 * 3} $\to$ \texttt{30} \\
\texttt{/} & деление & \texttt{10 / 3} $\to$ \texttt{3} \\
\texttt{\%} & остаток от деления & \texttt{10 \% 3} $\to$ \texttt{1} \\
\bottomrule
\end{tabular}
\end{table}

Деление двух целых даёт целое: дробная часть не округляется, а отбрасывается,
поэтому \texttt{10 / 3} равно 3, а не 3,333. Чтобы получить дробный результат,
дробным должен быть хотя бы один операнд: \texttt{10.0 / 3} или
\texttt{(double) 10 / 3}. Оператор \texttt{\%} возвращает остаток и чаще всего
служит проверкой делимости: \texttt{n \% 2 == 0}~--- признак чётного числа.

\begin{oshibka}
\textbf{Частая ошибка: целочисленное деление.} Многие ждут от \texttt{10 / 4}
значения 2,5, а получают 2, причём без единого предупреждения компилятора.
Особенно опасны формулы вида \texttt{c * 9 / 5 + 32}: если \texttt{c}~--- целое,
результат исказится. Лечится это одним символом: \texttt{c * 9.0 / 5 + 32}.
\end{oshibka}

Операторы сравнения \texttt{==}, \texttt{!=}, \texttt{>}, \texttt{<},
\texttt{>=}, \texttt{<=} возвращают \texttt{true} или \texttt{false}, а
логические \texttt{\&\&} («и»), \texttt{||} («или») и \texttt{!} («не»)
объединяют такие ответы. У первых двух есть важное свойство~--- \term{короткое
замыкание} (\eng{short-circuit}): если ответ ясен по левому операнду, правый не
вычисляется вовсе, и на этом строят защиту от деления на ноль. Увеличить
переменную на единицу тоже можно двумя способами, и разница видна, только когда
результат тут же используется: префиксный \texttt{++x} сначала прибавляет, а
потом отдаёт значение, постфиксный \texttt{x++}~--- наоборот. Оба приёма
показаны в листинге~\ref{lst:01-10}.

\begin{listing}[H]
\begin{minted}{java}
int age = 20;
boolean adult = age >= 18 && age < 65;  // true
System.out.println(!adult);             // false

int x = 0;
// Левый операнд уже ложен, поэтому 100 / x не вычисляется
if (x != 0 && 100 / x > 3) {
    System.out.println("не напечатается");
}

int n = 5;
System.out.println(++n);  // 6: сначала прибавили, потом взяли
System.out.println(n++);  // 6: сначала взяли, потом прибавили
System.out.println(n);    // 7
\end{minted}
\caption{Логические операторы, короткое замыкание и инкремент}
\label{lst:01-10}
\end{listing}

Так же устроены декременты \texttt{--x} и \texttt{x--}, а составные
присваивания \texttt{+=}, \texttt{-=}, \texttt{*=}, \texttt{/=}, \texttt{\%=}
сокращают запись \texttt{x = x + y} до \texttt{x += y}. Побитовые операторы
работают не с числом целиком, а с его двоичной записью, и нужны в работе с
флагами и масками: для \texttt{a = 5} (двоичное \texttt{0101}) и \texttt{b = 3}
(\texttt{0011}) побитовое И \texttt{a \& b} даёт 1 (\texttt{0001}), ИЛИ
\texttt{a | b}~--- 7 (\texttt{0111}), исключающее ИЛИ
\texttt{a \textasciicircum{} b}~--- 6 (\texttt{0110}), а сдвиги
\texttt{a << 1} и \texttt{a >> 1} умножают и делят на два, давая 10 и 2.

У оператора \texttt{+} есть и вторая работа: если хотя бы один его операнд~---
строка, остальные приводятся к тексту и склеиваются. Поэтому
\texttt{System.out.println("Вес: " + weight)} печатает подпись и число одной
командой, а \texttt{"1" + 2} даёт строку \texttt{"12"}, а не число 3.

\praktikum{}

Понадобится JDK~17 или новее и любой текстовый редактор. Проверьте установку
командами \texttt{java -version} и \texttt{javac -version}: обе должны
напечатать номер версии. Компилируйте с явным указанием кодировки~---
\texttt{javac~-encoding UTF-8 Имя.java}: иначе на JDK~17 под Windows русские
строки в исходнике будут прочитаны неверно (кодировкой по умолчанию UTF-8 стал
только в JDK~18). Если в консоли Windows вместо русского текста появляются знаки
вопроса, выполните перед запуском \texttt{chcp 65001}. Фрагменты кода из
теории~--- это тело метода: чтобы такой фрагмент запустить, поместите его внутрь
\texttt{public static void main(String[] args)} класса, имя которого совпадает с
именем файла, либо введите строки в \texttt{jshell}.

\subsection*{Задание 1. Первая программа и работа компилятора}

Создайте файл \texttt{HelloWorld.java} с программой, печатающей ваше имя и номер
группы. Скомпилируйте её командой
\texttt{javac~-encoding UTF-8 HelloWorld.java} и запустите командой
\texttt{java HelloWorld}. Запишите, какой файл появился после компиляции, затем
выполните \texttt{javap -c HelloWorld} и убедитесь, что видите байт-код.

Вторая часть~--- работа с ошибками. Наберите код из листинга~\ref{lst:01-11} в
файл \texttt{Errors.java} и скомпилируйте: в нём три ошибки. Исправляйте их по
одной, перезапуская компилятор, и выпишите текст каждого сообщения с
объяснением, что компилятору не понравилось.

\begin{listing}[H]
\begin{minted}{java}
public class Errors {
    public static void main(String[] args) {
        System.out.println("Привет, мир!")
        system.out.println("Вторая строка");
        System.out.println( Без кавычек );
    }
}
\end{minted}
\caption{Код с тремя ошибками для задания 1}
\label{lst:01-11}
\end{listing}

\subsection*{Задание 2. Эксперименты в jshell}

Запустите \texttt{jshell} и вычислите выражения \texttt{2 + 2},
\texttt{10 / 3}, \texttt{10.0 / 3}, \texttt{10 \% 3} и \texttt{Math.sqrt(16)}
(\texttt{Math}~--- класс из пакета \texttt{java.lang}, импорта не требует, а
\texttt{sqrt} возвращает \texttt{double}, поэтому ждите \texttt{4.0}, а не
\texttt{4}), каждый раз записывая прогноз до нажатия Enter.

Затем исследуйте границы типов. Объявите переменную \texttt{max} типа
\texttt{int} со значением \texttt{Integer.MAX\_VALUE} и вычислите
\texttt{max + 1}. Объявите \texttt{byte b = 127}, выполните \texttt{b++} и
посмотрите на значение \texttt{b}. Наконец, вычислите \texttt{0.1 + 0.2} и
\texttt{0.1 + 0.2 == 0.3}. Объясните письменно каждый из четырёх результатов.

\subsection*{Задание 3. Типы, литералы и ссылки}

Напишите программу \texttt{DataTypes.java}, в которой объявлены переменные всех
восьми примитивных типов и каждая напечатана на отдельной строке; не забудьте
суффиксы \texttt{L} для \texttt{long} и \texttt{f} для \texttt{float}.
Добавьте туда же четыре переменные типа \texttt{int} со значением 42
в десятичной, шестнадцатеричной, восьмеричной и двоичной записи и сравните их
оператором \texttt{==}, объяснив результат в отчёте.

Вторая часть~--- ссылки. В программе \texttt{StringComparison.java} объявите два
литерала \texttt{"Hello"} и два объекта, созданных через
\texttt{new String("Hello")}, и сравните все пары через \texttt{==} и
\texttt{equals}. Нарисуйте от руки схему памяти: покажите, что лежит в стеке, что
в куче и что в пуле строк, и стрелками отметьте, какая переменная на что
ссылается.

Третья часть~--- преобразования. Напечатайте значения \texttt{(int) 3.9},
\texttt{(byte) 300} и \texttt{(char)('A' + 1)}, а затем объясните, почему
\texttt{value += 300} для \texttt{byte} компилируется, а
\texttt{value = value + 300}~--- нет.

\subsection*{Задание 4. Небольшие расчётные программы}

Напишите две программы. Первая, \texttt{BMICalculator.java}, объявляет вес в
килограммах и рост в метрах как переменные типа \texttt{double}, вычисляет
индекс массы тела по формуле «вес, делённый на квадрат роста», и печатает
исходные данные и результат одной строкой, склеенной через \texttt{+}.
Вторая, \texttt{TemperatureConverter.java}, переводит температуру из шкалы
Цельсия в шкалу Фаренгейта по формуле $F = C \cdot 9 / 5 + 32$ и обратно по
формуле $C = (F - 32) \cdot 5 / 9$. Будьте внимательны с целочисленным делением:
проверьте, что дробная часть не теряется.

\subsection*{Результаты занятия}

По итогам занятия у вас должны быть готовы: файлы \texttt{.java} со всеми
заданиями и результатами запуска; объяснения к экспериментам из задания~2;
схема памяти из задания~3; сообщения компилятора из задания~1 с пояснением
каждой ошибки.

\kontrolvoprosy

\begin{enumerate}
\item Чем JDK отличается от JRE, а JRE~--- от JVM? Что нужно пользователю, а
      что~--- разработчику?
\item Что такое байт-код и какую задачу он решает?
\item Из каких трёх подсистем состоит JVM и за что отвечает каждая?
\item Перечислите три этапа загрузки класса. Что происходит на подэтапе
      проверки и зачем он нужен?
\item В чём состоит модель делегирования загрузчиков и какие три свойства
      системы она обеспечивает?
\item Чем куча отличается от стека и что хранится в каждой из этих областей?
\item Зачем нужен JIT-компилятор и что делает сборщик мусора?
\item Что означает слово \texttt{native} и когда за JNI браться не стоит?
\item Почему \texttt{value += 300} для переменной типа \texttt{byte}
      компилируется, а \texttt{value = value + 300}~--- нет?
\item Почему \texttt{10 / 3} даёт 3 и как получить дробный результат? В чём
      смысл короткого замыкания у операторов \texttt{\&\&} и \texttt{||}?
\end{enumerate}

\testzadaniya

\begin{enumerate}

\vopros{Во что компилятор \texttt{javac} превращает исходный код Java?}
  {в машинный код (\texttt{.exe})}
  {в байт-код (\texttt{.class})}
  {в ассемблерный код (\texttt{.asm})}
  {в текст на языке C}

\vopros{Какое соотношение между JDK, JRE и JVM верно?}
  {JVM содержит JRE, а JRE содержит JDK}
  {JRE содержит JDK, а JDK содержит JVM}
  {JDK и JRE~--- одно и то же}
  {JDK содержит JRE, а JRE содержит JVM}

\vopros{Что вернёт вызов \texttt{String.class.getClassLoader()}?}
  {Platform Class Loader}
  {Application Class Loader}
  {\texttt{null}}
  {строку \texttt{"Bootstrap"}}

\vopros{По какому принципу загрузчики классов ищут класс?}
  {каждый загрузчик ищет класс самостоятельно}
  {класс всегда загружает Application Class Loader}
  {поиск идёт одновременно сверху вниз и снизу вверх}
  {делегирование снизу вверх: запрос сначала уходит родителю}

\vopros{Где размещается объект, созданный оператором \texttt{new}?}
  {в стеке потока}
  {в Metaspace}
  {в куче}
  {в счётчике команд}

\vopros{Что делает JIT-компилятор?}
  {компилирует часто выполняемый байт-код в машинный код}
  {компилирует файлы \texttt{.java} в файлы \texttt{.class}}
  {загружает классы из JAR-архивов}
  {проверяет исходный код на ошибки перед запуском}

\vopros{Что такое Metaspace?}
  {область, в которой хранятся объекты}
  {стек вызовов методов}
  {область метаданных классов, заменившая PermGen}
  {кэш машинного кода, созданного JIT-компилятором}

\vopros{Каким ключевым словом объявляют метод, реализованный на C/C++?}
  {\texttt{extern}}
  {\texttt{foreign}}
  {\texttt{native}}
  {\texttt{external}}

\vopros{Каким будет результат \texttt{new String("Hi") == new String("Hi")}?}
  {\texttt{true}}
  {\texttt{false}}
  {\texttt{Hi}}
  {код не скомпилируется}

\vopros{Что произойдёт при выполнении \texttt{byte b = 127; b++;}?}
  {код не скомпилируется}
  {переменная \texttt{b} станет равна 128}
  {будет выброшено исключение}
  {переменная \texttt{b} станет равна $-128$ из-за переполнения}

\vopros{Почему выражение \texttt{0.1 + 0.2 == 0.3} возвращает \texttt{false}?}
  {оператор \texttt{==} не работает с типом \texttt{double}}
  {нужно было использовать \texttt{float} вместо \texttt{double}}
  {из-за погрешности двоичного формата IEEE~754}
  {Java округляет дробные числа при сложении}

\vopros{Что напечатает \texttt{System.out.println(10 / 3)}?}
  {\texttt{3.33}}
  {\texttt{3}}
  {\texttt{3.0}}
  {\texttt{4}}

\vopros{Что произойдёт при вычислении \texttt{false \&\& someMethod()}?}
  {\texttt{someMethod()} будет вызван, результат~--- \texttt{false}}
  {\texttt{someMethod()} не будет вызван: короткое замыкание}
  {произойдёт ошибка компиляции}
  {результат зависит от того, что вернёт \texttt{someMethod()}}

\vopros{Что означает литерал \texttt{0b101010}?}
  {восьмеричное число}
  {шестнадцатеричное число}
  {двоичную запись числа 42}
  {строку из нулей и единиц}

\end{enumerate}
